\section{Conclusion}

\label{sec:conclusion}
%% ============================================================================

Quasar is the first production consensus protocol to require three independent
cryptographic assumptions for finality. Safety holds if at least one layer
remains secure. Liveness is guaranteed under partial synchrony with $f < n/3$
Byzantine validators. The protocol is implemented in Go, deployed on Lux
Network, and supported by machine-checked proofs in Lean~4 and symbolic
verification in Tamarin.

Performance is practical: $69\,\text{ms}$ full threshold rounds at $n=100$
validators, $18.8\,\text{ns}$ finality verification, and $20.26$M messages
per second on the ZAP transport layer. Epoch-based pruning bounds memory at
$5.9\,\text{MB}$ regardless of chain length. The test suite covers $2{,}160+$
assertions across all consensus modes.

The quantum-layer design is deliberately conservative. We do not claim that
BLS12-381 will resist quantum computers, nor that Ring-LWE or Module-LWE will
survive the next two decades of cryptanalysis. We claim that requiring an
adversary to break all three simultaneously---across three algebraically
independent hardness assumptions---provides a margin of safety that no
single-scheme approach can match.

\bibliographystyle{plain}
\begin{thebibliography}{99}

\bibitem{shor1994}
P.~W.~Shor.
\newblock Algorithms for Quantum Computation: Discrete Logarithms and
  Factoring.
\newblock In \emph{FOCS}, 1994.

\bibitem{nist-pqc}
{National Institute of Standards and Technology}.
\newblock Post-Quantum Cryptography Standardization.
\newblock NIST, 2024.

\bibitem{mosca2018}
M.~Mosca.
\newblock Cybersecurity in an Era with Quantum Computers.
\newblock \emph{IEEE Security \& Privacy}, 2018.

\bibitem{boneh2001short}
D.~Boneh, B.~Lynn, and H.~Shacham.
\newblock Short Signatures from the Weil Pairing.
\newblock In \emph{ASIACRYPT}, 2001.

\bibitem{ringtail}
Z.~Kelling and V.~Seesahai.
\newblock Ringtail: Lattice-Based Post-Quantum Threshold Signatures for
  Blockchain Consensus.
\newblock Lux Industries, 2024.

\bibitem{castro1999practical}
M.~Castro and B.~Liskov.
\newblock Practical Byzantine Fault Tolerance.
\newblock In \emph{OSDI}, 1999.

\bibitem{yin2019hotstuff}
M.~Yin, D.~Malkhi, M.~K.~Reiter, G.~G.~Gueta, and I.~Abraham.
\newblock HotStuff: BFT Consensus with Linearity and Responsiveness.
\newblock In \emph{PODC}, 2019.

\bibitem{groth2016}
J.~Groth.
\newblock On the Size of Pairing-Based Non-Interactive Arguments.
\newblock In \emph{EUROCRYPT}, 2016.

\bibitem{regev2005}
O.~Regev.
\newblock On Lattices, Learning with Errors, Random Linear Codes, and
  Cryptography.
\newblock \emph{J.~ACM}, 56(6), 2009.

\bibitem{espitau2024}
T.~Espitau, M.~Tibouchi, and A.~Wallet.
\newblock Practical Fiat-Shamir with Aborts Threshold Signatures from
  Lattices.
\newblock \emph{IACR ePrint}, 2024.

\bibitem{falcon-spec}
P.-A.~Fouque, J.~Hoffstein, P.~Kirchner, V.~Lyubashevsky, T.~Pornin,
  T.~Prest, T.~Ricosset, G.~Seiler, W.~Whyte, and Z.~Zhang.
\newblock FALCON: Fast-Fourier Lattice-Based Compact Signatures over NTRU.
\newblock NIST PQC Round~3 submission, 2020.

\bibitem{ducas2016}
L.~Ducas and T.~Prest.
\newblock Fast Fourier Orthogonalization.
\newblock In \emph{ISSAC}, 2016.

\bibitem{qcert-proof}
Z.~Kelling.
\newblock QuasarCert Soundness: Proof Sketch with Circuit Analysis,
  Adaptive Corruption, Setup, and Parameter Tightness.
\newblock Lux Industries technical note, April 2026.
\newblock \texttt{proofs/quasar-cert-soundness.tex}.

\end{thebibliography}

%% ============================================================================
\appendix
%% ============================================================================

